\documentclass[12pt]{article}

\input{.house-style/preamble.tex}
\usepackage{xurl}
\usepackage{booktabs}
\usepackage{array}
\setlength{\emergencystretch}{3em}
\setlength{\headheight}{14pt}

\addbibresource{references-local.bib}

% Fix duplicate DOI on ameka1992 and add missing URL on davies2020coca.
\DeclareSourcemap{
  \maps[datatype=bibtex]{
    \map[overwrite=true]{
      \step[fieldsource=entrykey, match=\regexp{\Aameka1992\Z}, final]
      \step[fieldset=doi, fieldvalue={10.1016/0378-2166(92)90048-G}]
    }
    \map[overwrite=true]{
      \step[fieldsource=entrykey, match=\regexp{\Adavies2020coca\Z}, final]
      \step[fieldset=url, fieldvalue={https://www.english-corpora.org/coca/}]
    }
  }
}

\newif\ifblind
\blindfalse
\ifblind
\DeclareSourcemap{
  \maps[datatype=bibtex]{
    \map[overwrite=true]{
      \step[fieldsource=entrykey, match=\regexp{\Areynolds}, final]
      \step[fieldset=author, fieldvalue={Author}]
      \step[fieldset=title, fieldvalue={[Title removed for anonymous review]}]
      \step[fieldset=note, null]
      \step[fieldset=url, null]
      \step[fieldset=doi, null]
    }
  }
}
\fi

% Breakable monospace macro preserving visible token spaces in query tables.
\newcommand{\qsep}{\textbar\allowbreak}

\hypersetup{
  pdftitle={Damn him to hell, fuck me for giving a shit, thank you for coming: How Far Generalizations Reach in English},
  pdfkeywords={imprecatives, benedictives, construction grammar, defective verbs, English syntax},
}
\ifblind
\hypersetup{pdfauthor={}}
\else
\hypersetup{pdfauthor={Brett Reynolds}}
\fi

\title{\mention{Damn him to hell}, \mention{fuck me for giving a shit}, \mention{thank you for coming}:\\How Far Generalizations Reach in English}
\ifblind
\author{[Author removed for blind review]}
\else
\author{Brett Reynolds \orcidlink{0000-0003-0073-7195}%
\thanks{Contact: \href{mailto:brett.reynolds@humber.ca}{brett.reynolds@humber.ca}}\\
Humber Polytechnic \& University of Toronto}
\fi
\date{\today}

\begin{document}
\maketitle

\begin{abstract}
Grammatical accounts assign different categories to English expressions without an overt subject, such as \mention{Damn the bell!} and \mention{Bless you!}. This paper asks which properties recur across the sampled heads and which identify narrower or wider groupings. I compare thirteen heads, checking by hand every candidate retrieved from 20 million words. Searches of 12 million words of unscripted speech and 2.2 billion of subtitles and web text extend the comparison to six \mention{fuck}-euphemisms and \mention{confound}. \mention{To hell} concentrates on \mention{damn} (32 of 33 audited shared-pronoun tokens), overt \mention{God} on \mention{damn} and particular \mention{bless} formulae, and evaluative \mention{for} V-ing occurs with six heads and with \mention{thank you}. Open NPs recur across most heads; euphemisms share many of their sources' frames in different proportions. Non-reflexive \mention{you} and \mention{me} exclude a uniform addressee or speaker subject, and self-curses contradict a semantic exclusion of first-person targets. The TV Corpus has \mention{damn it to hell} but no \mention{damn it for} V-ing. Each generalization has a different scope: extending \mention{damn}'s or \mention{fuck}'s profile to the other heads makes wrong predictions, and no property studied picks out exactly the inherited set. The evidence constrains subjects and extensions while leaving lexical category and grammatical representation open.
\end{abstract}

\noindent\textbf{Keywords:} imprecatives; benedictives; construction grammar; defective verbs; English syntax

% ===========================================================================
\section{Introduction}
\label{sec:intro}
% ===========================================================================

English has a pattern in which an expressive head combines directly with a noun phrase and, optionally, a prepositional phrase, without an overt grammatical subject. Grammatical accounts disagree about what category its heads belong to; this paper asks which of its properties hold across the heads those accounts argue about and which belong to narrower or wider groupings, and so what any one category, lexical or otherwise, can predict. Attested examples span secular taboo maledictives, theological curses, affectionate benedictives, and mild minced oaths, as illustrated in (\ref{ex:intro-paradigm}).

\ea\label{ex:intro-paradigm}
\ea \mention{Damn the bell!} \hfill (BNC Baby \texttt{H85})
\ex \mention{Fuck the doctors.} \hfill (BNC Baby \texttt{CCW})
\ex \mention{Sod the choir!} \hfill (BNC Baby \texttt{KPU})
\ex \mention{Bugger the rest!} \hfill (CABNC \texttt{KB7})
\ex \mention{Bless you for helping me out with the Watsons.} \hfill (BNC Baby \texttt{G0Y})
\ex \mention{Oh, darn the luck.} \hfill (TV Corpus, \emph{The Simpsons})
\z
\z

Since \textcite[3--7]{quang1971} analysed \mention{fuck you} and \mention{damn him} as \enquote{English sentences without overt grammatical subject} headed by a distinct category of \enquote{quasi-verbs}, grammarians have struggled to locate this pattern in English syntax. \textcite[236, 944, 1361 n.\ 70]{huddleston2002} distinguish restricted verbal optatives from \mention{Damn these mosquitoes!}, tentatively treating the latter as an exceptional interjection phrase and doubting an understood subject. \textcite[852--853]{quirk1985} list \mention{Damn} (\mention{it}), \mention{Bugger} (\mention{it}), and \mention{Fuck} (\mention{it}) among expletive formulae alongside interjections.

Semanticists and corpus linguists approach the same forms from different angles. Formal semanticists model expressive and use-conditional meaning \citep{potts2007,gutzmann2013,gutzmann2015,gutzmann2019}, and \textcite[4]{portner2026} gives profane quasi-verbs a social-relation semantics while setting their morphosyntax aside. Section~\ref{sec:meaning} asks what the distributions reported here require of such accounts. Corpus studies treat them as taboo verbs, euphemisms, or interactives \citep{mcenery2006,allanburridge2006,napolihoeksema2009,ljung2011,heine2023interactives}, and \textcite[13]{ljung2011}, discussing their shared swearing function, calls \mention{fuck}, \mention{damn} and \mention{sod} \enquote{freely interchangeable in curses like \mention{Fuck you!}, \mention{Damn you!} and \mention{Sod you!}}.

The choice among analyses (a defective verb, an interjection, a slot in a construction) raises four separable questions: (1) separation from primary interjections, \mention{yay}/\mention{boo}, imperatives, and uniform reduced performatives; (2) which extensions (\mention{God}, \mention{to hell}, \mention{for} V-ing, the bare-subjunctive \mention{be} + participle frame, coordination) generalize vs.\ stay lexically anchored; (3) complement-NP openness vs.\ which words fill the head slot, including partial euphemistic correspondence; and (4) what these uneven distributions mean for categorizing these expressions.

The terms \term{imprecative} and \term{benedictive} name pragmatic categories, what a speaker does, whereas the pattern in (\ref{ex:intro-paradigm}) is a form, and the two don't coincide. Curses and blessings take many other forms: the \mention{may}-optative (\mention{May you prosper!}), the imperative (\mention{Be well!}, \mention{Go to hell!}; \citealp[32]{ljung2011}), the subjunctive with an overt subject (\mention{God damn you}), preposition-headed dismissals (\mention{to hell with} X), finite clauses with a modal (\mention{\dots\ can fuck himself}; \texttt{D9}) and declaratives (\mention{You're a dirty rotten scoundrel and you deserve to die}). None combines a bare head with an NP complement and no overt subject, so what this paper finds about the pattern isn't a finding about cursing or blessing in general.

The same form also serves acts that aren't curses or blessings: thanks (\mention{Thank you}), dismissal (\mention{Hang the expense}, TV Corpus, \emph{Cold Feet}) and surprise (\mention{Fuck me. That's incredible news}, TV Corpus, \emph{White Gold}). The configuration alone admits ordinary imperatives too (\mention{Close the door!}), so the object of study is the expressive uses of the sampled heads within it, picked out by form, by a list of heads and by use together. This study samples the pattern through the thirteen heads that grammatical accounts have argued about (Section~\ref{sec:diagnostics}), six euphemisms of \mention{fuck} and archaic \mention{confound}. Those heads are a sampling frame inherited from the literature, not a closed inventory (Section~\ref{sec:tier1-regional}), and \mention{thank} is the nearest comparison outside it (Section~\ref{sec:beyond}).

This study addresses the four questions by checking by hand every candidate retrieved from 20M words (CABNC, BNC Baby XML from the British National Corpus, BNC, and the Open American National Corpus, OANC) and by searching 12M words of unscripted British speech (Spoken BNC2014) and 2.2B words of subtitles and web text (TV Corpus, GloWbE), with a few exploratory checks in the Corpus of Contemporary American English (COCA). Its aim is comparative grammar with explicit bounds on incomplete displays: mapping what's attested, what's lexically concentrated, and what observational data do and don't settle. Section~\ref{sec:convergence} sorts the results by the scope of each generalization.

% ===========================================================================
\section{Grammatical profile, pragmatic subtypes, and analytical positions}
\label{sec:frameworks}
% ===========================================================================

\subsection{Morphosyntactic diagnostics}
\label{sec:diagnostics}

The thirteen core heads come from three historical strata: theological curses and blessings (\mention{damn}, \mention{bless}, \mention{curse}, \mention{blast}), secular taboo or violent verbs (\mention{fuck}, \mention{screw}, \mention{sod}, \mention{bugger}, \mention{hang}), and minced oaths (\mention{drat}, \mention{darn}, \mention{dash}, \mention{dang}).\footnote{Historically, \mention{drat} contracts (\mention{Go})\mention{d rot}, \mention{darn} and \mention{dang} are phonological euphemisms for \mention{damn}, and \mention{dash} entered early nineteenth-century British English as a mild euphemism associated with the printed dash replacing \mention{damn} \citep{mcenery2006,allanburridge2006,ljung2011}.} Six of the heads, \mention{damn}, \mention{bless}, \mention{fuck}, \mention{screw}, \mention{sod} and \mention{bugger}, are in everyday contemporary use, and I call these the \term{living heads} (Section~\ref{sec:tv-profile}). Across the strata, the pattern shows the seven grammatical properties summarized in (\ref{ex:diagnostics}), though not every head shows every one (Section~\ref{sec:convergence}); asterisks mark ungrammaticality \emph{on the expressive reading}. Here and below, \term{subjectless} means without an overt subject; whether there's an understood one is what Positions A--D (Section~\ref{sec:four-positions}) disagree about.

\ea\label{ex:diagnostics}
\ea Obligatory accusative pronoun form on case-marking pronouns: \mention{Damn him!} / \mention{Fuck them!} / \ungram{\mention{Damn he!}} / \ungram{\mention{Fuck they!}}
\ex Non-reflexive second- and first-person complements: \mention{Fuck you!} / \mention{Damn you all!} (vs.\ imperative \mention{Fuck yourself!}); \mention{Damn me for agreeing to it!} / \mention{Fuck me!} (vs.\ \ungram{\mention{I damn me for agreeing to it!}}, predicting \mention{myself})
\ex Open lexical NP complementation: \mention{Screw the Times!} / \mention{Sod the bar codes!} / \mention{Oh, darn the luck!} / \mention{Confound that man!}
\ex Absence of tense, agreement, aspect, imperative \mention{do}-negation, and \mention{please} (on the expressive reading): \ungram{\mention{Damned the bell!}} / \ungram{\mention{Don't damn him!}} (imprec.) / \odd{\mention{Please sod the choir!}}
\ex Absence of ordinary overt subjects (on the expressive reading): \ungram{\mention{I sod the choir!}} / \ungram{\mention{You fuck you!}} (vs.\ imperative \mention{You fuck off!})
\ex No eventive passive preserving imprecative force (\ungram{\mention{The rules are being fucked by me!}}, imprec.), alongside bare-subjunctive \mention{be} + participle frames on some heads: dismissals (\mention{The law be damned!}, \mention{Sony be buggered!}) and, inverted, benedictions (\mention{Blessed be the name of the Lord})
\ex Split modification: rejection of event-level manner and time adjuncts on \mention{fuck}/\mention{damn} (\ungram{\mention{Fuck him quickly!}}, \ungram{\mention{Damn those irregular verbs tomorrow!}}; \citealp[5--6]{quang1971}) alongside evaluative \mention{for} V-ing (\mention{Fuck her for killing herself!}), dismissal-scope temporal adjuncts (\mention{Screw the Times \dots\ for the moment}), and lexically restricted \mention{to hell} (\mention{Damn him to hell!}).
\z
\z

Three points sharpen (\ref{ex:diagnostics}). First, accusative form in (\ref{ex:diagnostics}a) doesn't alone prove verbal Case assignment: English assigns \term{default accusative case} outside nominative positions, including non-verbal expressive small clauses (\mention{Poor me!}, \mention{Silly me!}; \citealp{schutze2001,pottsroeper2006}) and expressive \mention{yay me!} / \mention{boo him!} (Section~\ref{sec:beyond}). In these corpora, what separates (\ref{ex:diagnostics}) from primary interjections (\mention{ouch}, \mention{oops}, \mention{ugh}, \mention{wow}, \mention{whoops}, \mention{yuck}; \citealp{ameka1992,amekawilkins2006,meinard2015}) and from expressive \mention{yay}/\mention{boo} (Sections~\ref{sec:tv-profile} and~\ref{sec:bnc2014}) is open lexical NPs. Outside these corpora, \mention{yay} also occurs with \mention{for} V-ing and with a lexical NP (\mention{Yay the NHS!}; Section~\ref{sec:beyond}).

Second, non-reflexive complements in (\ref{ex:diagnostics}b) test for an understood subject. An addressee subject would bind a reflexive, as in imperative \mention{Fuck yourself!} \citep[3]{quang1971,pinker2007,portner2026}, so non-reflexive \mention{you} excludes one for the \mention{you} forms. All 56 inspected TV lines of \mention{fuck myself} have a subject or controller, or echo a previous turn; none is a self-curse (Supplementary Section~S5).

A speaker subject, as in reduced-performative \mention{Thank you} $<$ \mention{I thank you}, would require \mention{myself} when the speaker is the target. Non-reflexive \mention{me} in self-curses (\mention{Damn me for agreeing to it}, \emph{Babylon 5}; \mention{Fuck me for giving a shit}, \emph{Shameless}; TV Corpus) excludes one for those forms. Neither test excludes a third-person subject, and an addressee subject remains possible for \mention{me} with someone present (Section~\ref{sec:ruled-out}). Literal readings remain grammatical in (\ref{ex:diagnostics}d--f); \odd{\mention{Please sod the choir!}} reflects register clash rather than pure syntactic ill-formedness.

Third, (\ref{ex:diagnostics}f--g) qualify two classical claims. The bare-subjunctive \mention{be} + participle frames in (\ref{ex:diagnostics}f) show that some heads also occur as participles outside the active frame (Section~\ref{sec:passive}). And while \textcite[5--6]{quang1971} cited the rejection of event-level manner and time adjuncts (\ungram{\mention{Damn those irregular verbs tomorrow!}}) against verbhood, Sections~\ref{sec:tier1} and \ref{sec:dissoc} document productive modifiers that scope over the expressive act, not an event, especially \mention{for} V-ing and dismissal-scope \mention{for the moment}. I call a reason PP such as \mention{for} V-ing \term{evaluative} when it states the grounds of an attitude the expression directs at a target: in \mention{Fuck her for killing herself!}, the grounds of the curse aimed at her (Section~\ref{sec:meaning} relates the notion to semantic accounts).

The diagnostics in (\ref{ex:diagnostics}) differ in standing. (\ref{ex:diagnostics}d--e) partly restate how tokens were admitted, since the large-corpus queries search bare forms and a token with an overt subject was coded as an ordinary verb. \texttt{C4} gives an independent check for \mention{damned} and \mention{cursed}: all 100 sampled GloWbE lines with a pronoun subject are ordinary finite verbs. The rejection of event-level adjuncts in (\ref{ex:diagnostics}g) rests on judgments \citep[5--6]{quang1971}.

(\ref{ex:diagnostics}a) is attested in every search, but only the BNC Baby search, which admits any pronoun, could have found a nominative. The non-reflexive complements in (\ref{ex:diagnostics}b), the open NPs in (\ref{ex:diagnostics}c), the \mention{be} + participle frames in (\ref{ex:diagnostics}f), \mention{for} V-ing and \mention{to hell} are findings.

\subsection{Pragmatic and semantic subtypes}
\label{sec:subtypes}

Uses of HEAD + NP (HEAD standing for any sampled head) differ pragmatically in ways that bear on their extensions. Dismissals often target inanimate things or institutions (\mention{Damn the bell!}, \mention{Screw the committee}), whereas targeted curses and acts of praise or thanks can evaluate a referent for something (\mention{Fuck her for killing herself}, \mention{Bless him for trying}). The reason PP states the grounds for that attitude. It also helped identify these uses, so its distribution doesn't independently test a requirement that the target be held responsible (Section~\ref{sec:dissoc}). Routine uses vary in how much reference they keep: formulaic \mention{Damn it!} has non-referential \mention{it}, whereas conventionalized \mention{Bless you!} bleaches the blessing act while still picking out its recipient through \mention{you}. To my ear, \mention{Bless it!} needs a referent (a baby or a pet, say), where \mention{Damn it!} and \mention{Fuck it!} need none. The corpora are consistent with that judgment without testing it: Tier~1's only \mention{bless it} is referential (Supplementary Table~S3, note a), and the TV Corpus has 51 raw \mention{bless it} strings against 20,107 for \mention{damn it} (Supplementary Table~S4).

\subsection{Live analytical accounts}
\label{sec:four-positions}

Four analytical perspectives frame the literature. Under optative and reduced-performative subject omission (Position~A), the forms have an understood subject. For \mention{bless} and \mention{damn} it's a third-person subject referring to God, motivated by their history as \mention{God}-optatives \citep{huddleston2002,ljung2011}; in reduced performatives (\mention{Thank you} $<$ \mention{I thank you}) it's the speaker, which would account for non-reflexive \mention{you}. The \mention{for} phrase in \mention{thank you for coming} supplies a comparison for reason modification (Section~\ref{sec:beyond}).

A second perspective contrasts primary interjections with complement-retaining interjection phrases (Position~B). Interjection theories separate \term{primary interjections} (\mention{ouch}, \mention{oops}, \mention{ugh}), which don't govern NP complements, from \term{secondary interjections} (\mention{damn!}, \mention{fuck!}, \mention{bless!}) recruited from lexical words \citep{ameka1992,amekawilkins2006,meinard2015,heine2023interactives}. On this view, the heads studied here head complement-retaining interjection phrases \citep[1361 n.\ 70]{huddleston2002}.

A third perspective treats the heads as defective verbs (Position~C). In a lexicalist framework \citep{mullerwechsler2014}, each head could carry a defective verbal entry selecting an NP in a subjectless frame. Missing finite forms would follow from the entry's restricted paradigm (compare \mention{beware}, \mention{betide}). Quang's separate quasi-verb category is another lexical analysis, rather than a defective-verb account \citep[5--7]{quang1971}.

A fourth perspective is a usage-based constructional network (Position~D). Drawing on modern Construction Grammar \citep{fillmore1988mechanisms,goldberg1995constructions,goldberg2003,croft2001,bybee2010,diessel2019,goldberg2019,schmid2020}, one can posit a structured network pairing schematic [X\textsubscript{expr} NP] with lexically anchored sub-constructions (civic \mention{God bless} NP, \mention{damn} NP \mention{to hell}, British \mention{damn and blast}) and horizontal euphemistic links.

The four positions answer different questions: whether there's an understood subject and what it could refer to (Position~A, and the binding tests of Section~\ref{sec:diagnostics}); what lexical category the head belongs to (Positions~B and~C); how head-specific extensions are organized (all four); and where that knowledge is represented (Positions~C and~D). A further question, what the expressions mean and how their dependents combine with that meaning, belongs to the semantic accounts cited in Section~\ref{sec:intro}, which Section~\ref{sec:meaning} takes up. Some of their commitments can coexist, so the evidence below constrains different positions on different questions.

% ===========================================================================
\section{Data and methods}
\label{sec:methods}
% ===========================================================================

No single corpus pairs exhaustive manual checking with the scale that rare frames need, so I combine two tiers.

\subsection{Tier 1: British and American corpora checked in full (20M words)}
\label{sec:tier1-methods}

Tier~1 searches three corpora in full: CABNC \citep{albert2015cabnc} (2.05M tokens of British conversation from BNC1994), BNC Baby XML \citep{burnard2007bncbaby} (4.0M tokens, CLAWS5-tagged), and OANC \citep{ide2008oanc} (14.45M tokens of American English). Removing 12 overlapping BNC Baby files leaves 19,947,325 distinct tokens (5.50M UK; 14.45M US).

The CABNC and OANC searches require the NP to begin with a closed-class word (\mention{it}, \mention{you} or an accusative pronoun, a determinative, a genitive pronoun, or a quantifier such as \mention{all} or \mention{everyone}). The BNC Baby search accepts any NP-initial tag, as a check on what the closed-class searches miss. I checked every one of the 596 candidates by hand, along with 148 candidates for six primary interjections (\mention{ouch}, \mention{oops}, \mention{ugh}, \mention{wow}, \mention{whoops}, \mention{yuck}). CABNC transcripts carry no timed pauses (none in 200 files inspected), so the BNC1994 transcribers' sentence-initial capitals (\mention{Ouch My foot}) serve only as a rough cue to unit boundaries.

\subsection{Tier 2: Subtitle, web, and contemporary spoken corpora}
\label{sec:tier2-methods}

Tier~2 adds scale. The TV Corpus \citep{davies2019tv} holds 325M words of subtitles (1950--2018) and GloWbE \citep{davies2013glowbe,daviesfuchs2015} 1.9B words of web text from 20 varieties; a few exploratory checks use COCA \citep{davies2020coca}. Spoken BNC2014 \citep{loveEtAl2017SpokenBNC2014} (11.83M words of unscripted British conversation, 2012--2016) tests whether the subtitle distributions hold in speech. Supplementary Table~S1 lists every query with the ID used in the tables (\texttt{T7a}, for example, is HEAD + pronoun + \mention{to hell} in the TV Corpus), and the replication archive gives full tallies for each.

The TV Corpus interface displays at most 100 concordance lines per query by default, and not a random 100. Where a query returns more hits than it displays, I audit the displayed lines and report bounds that hold however the undisplayed lines turn out (Section~\ref{sec:dissoc}). The Chart view, by contrast, gives a sample of up to 100 lines drawn from all of a query's hits (every line, for smaller queries). I use it for GloWbE and for eight TV searches (five of the frames in Supplementary Table~S4, \mention{confound}, \mention{be darned} and the euphemisms of \mention{fuck}), and audit the samples in the same way.

\subsection{Coding and quantitative reporting}
\label{sec:coding}

Tier-1 coding distinguishes \mention{it}, routine \mention{bless you}, other pronouns, demonstratives, body-part and soul formulae, and open full NPs; overt-\mention{God} optatives are counted separately. Supplementary Table~S3 gives the complete inventory and Section~S3 the coding priorities and exclusions.

Exclusions remove ordinary transitive verbs (\mention{to hang the curtains}, \mention{to darn her own stockings}), negative quantifier idioms (\mention{bugger/sod/fuck all}), standalone interjections before a new clause (\mention{Fuck. That would have been amusing}), citations of the word, proper nouns (\emph{Dash}), and West African ditransitive \mention{dash} (\enquote{give, tip}) in GloWbE.

Tokens cluster by file, formula and series, and the displayed concordance lines aren't random samples, so I report no \emph{p}-values \citep{kilgarriff2005}. Counts retain their denominators and dispersion across files, shows or conversations. A zero counts as a nonattestation unless a rate from another head or corpus gives a benchmark for how many tokens to expect. Supplementary Sections~S6 and~S9 give the rate calculations and sensitivity checks.

% ===========================================================================
\section{Tier-1 results}
\label{sec:tier1}
% ===========================================================================

Tier~1 yields 596 candidates, of which 220 are verified subjectless tokens from 143 files; 23 are optatives with overt \mention{God}, and 353 are excluded (Supplementary Table~S3). Nine of the thirteen heads are attested, and six take full lexical NPs (24 tokens: \mention{damn} 7, \mention{fuck} 5, \mention{screw} 5, \mention{sod} 4, \mention{bugger} 2, \mention{bless} 1).


Tier~1 attests the modifiers that (\ref{ex:diagnostics}g) allows. The \mention{for}-phrases in (\ref{ex:tier1-gems}) modify the phrase headed by \mention{damn}, \mention{fuck}, \mention{sod} or \mention{screw}, and they're of three kinds: the reason for a curse (\ref{ex:tier1-gems}a), British idiomatic dismissals (\ref{ex:tier1-gems}b--c), and a time adjunct, \mention{for the moment}, that scopes over the act of dismissing rather than over an event inside it (\ref{ex:tier1-gems}d).

\ea\label{ex:tier1-gems}
\ea \mention{Damn the myopic but all-seeing eyes of that news-ferret Cramer for having the cruel wit to reduce him to an onomatopoeia.} \hfill (OANC \texttt{Slate})
\ex \mention{Oh fuck this for an idea!} \hfill (CABNC \texttt{KCERE005})
\ex \mention{\dots would be to actually say sod this for a lark\dots} \hfill (CABNC \texttt{KBHRE005})
\ex \mention{Screw the Times, young readers, and everything else for the moment.} \hfill (OANC \texttt{Slate})
\z
\z

\subsection{No NP complements in the six controls}
\label{sec:cabnc-boundaries}

The six primary interjections give 148 candidates followed by an NP-initial word, and none is a direct NP complement. Every hit is the subject of a following clause (\mention{Whoops, you've dropped your pencil}), a nursery formula (\mention{oops a daisy}), a lexical verb (\mention{to wow us}), or a detached fragment (\mention{Ouch. My foot}; \mention{Ugh, the prep}; compare \ungram{\mention{Ouch him!}}, \ungram{\mention{Ugh them!}}).

Transcriber capitalization supports this distinction but partly entered the coding, so it supplies no independent prosodic test (Supplementary Section~S3).

\subsection{Missed heads, regional skew, and overt \mentionhead{God}}
\label{sec:tier1-regional}

The broader BNC Baby search recovered three bare nouns or proper names that the closed-class searches miss (\mention{Damn McIllvanney}, \mention{Oh, screw Jesus Christ}, \mention{night, bless mummy and dad}). It also exposed an ambiguity: \mention{Damn DTs.} is either a subjectless imprecative or prenominal \mention{damn} (compare \mention{Damn rubbish.}; \citealp[35 n.\ 29]{portner2026}). That ambiguity is why the large-corpus queries require an overt determinative, genitive pronoun or accusative pronoun.

The British data strongly favours \mention{sod} and \mention{bugger}, but register differences prevent a controlled UK--US comparison. Overt \mention{God} concentrates on \mention{damn} and \mention{bless} (Supplementary Section~S3), a distribution tested at scale below.

A broader search also finds \mention{blow} (\mention{blow this for a lark}), \mention{stuff} (\mention{stuff it!}) and \mention{stone} (\mention{stone the crows}; Supplementary Section~S3). The thirteen heads are a sampling frame, not a closed inventory. Outside it, \mention{Thank you} shares the second-person binding facts, though it hasn't been shown to share the first-person one, and \mention{Fancy that!} expresses neither a curse nor a blessing. Section~\ref{sec:beyond} compares \mention{thank}.

% ===========================================================================
\section{Tier-2 results}
\label{sec:tier2}
% ===========================================================================

\subsection{The TV Corpus profile and its controls}
\label{sec:tv-profile}

In 100-line TV samples, subjectless imprecatives and benedictives account for 78 of the HEAD + \mention{it} lines but only 35 of the HEAD + genitive + N lines. Raw strings include ordinary verbs, overt-subject optatives and other uses. Supplementary Table~S4 and Section~S4 give the complete inventory and audits.

The concordances verify open lexical NPs across the living heads (\mention{Screw the committee, screw the party}, \emph{The Art of More}; \mention{Sod the bar codes!}, \emph{Code of a Killer}), and beyond them (\mention{Hang the casting committee}, \emph{Room at the Top}; \mention{Oh, darn the luck}, \emph{The Simpsons}). None of 53 primary-interjection + pronoun lines or 13 \mention{yay}/\mention{boo} + lexical-NP candidates has an expressive complement.

Coordination of two heads over one pronoun gives 11 \mention{damn and blast} tokens, all in British period drama or comedy; speakers also coordinate whole imprecatives (\mention{Fuck the Store and fuck the Troubadour}). Archaic \mention{dash} and \mention{confound} are concentrated in period drama: one series supplies 43 of 75 audited \mention{confound} imprecatives.

\subsection{\mentionhead{To hell} versus \mentionhead{for} V-ing}
\label{sec:dissoc}

The two PP frames split the heads. Table~\ref{tab:t7-reconciled} compares HEAD + pronoun + \mention{to hell} with HEAD + pronoun + \mention{for} V-ing. The two queries used slightly different pronoun lists, so the comparison rests on the pronouns both share (\mention{you/him/her/them}), which leaves out bleached \mention{damn it to hell} and \mention{me}.

\begin{table}[tbp]
\centering
\caption{Audited subjectless PP tokens on shared pronouns (\mention{you/him/her/them}) in the TV Corpus. Bases are raw head + pronoun strings, including ordinary verbs and overt-subject optatives; token counts exclude those uses. Full query totals and exclusions are in Supplementary Table~S5.}
\label{tab:t7-reconciled}
\footnotesize
\setlength{\tabcolsep}{6pt}
\begin{tabular}{@{}lrrrrr@{}}
\toprule
 & Raw base & \multicolumn{2}{c}{\mention{to hell}} & \multicolumn{2}{c}{\mention{for} V-ing} \\
\cmidrule(lr){3-4}\cmidrule(lr){5-6}
Head & & Tokens & Shows & Tokens & Shows \\
\midrule
\mention{damn} & 1,627 & 32 & 29 & 10 & 10 \\
\mention{bless} & 3,113 & 0 & 0 & 6 & 6 \\
\mention{fuck} & 8,202 & 0 & 0 & 17 & 17 \\
\mention{screw} & 2,593 & 0 & 0 & 13 & 12 \\
\mention{curse} & 298 & 0 & 0 & 1 & 1 \\
\mention{blast} & 357 & 1 & 1 & 0 & 0 \\
\mention{darn} & 33 & 0 & 0 & 1 & 1 \\
\mention{dang} & 5 & 0 & 0 & 0 & 0 \\
\mention{sod} & 51 & 0 & 0 & 0 & 0 \\
\mention{bugger} & 24 & 0 & 0 & 0 & 0 \\
\mention{drat} & 6 & 0 & 0 & 0 & 0 \\
\mention{dash} & 14 & 0 & 0 & 0 & 0 \\
\mention{hang} & 1,008 & 0 & 0 & 0 & 0 \\
\midrule
Total & 17,331 & 33 & 29 & 48 & 47 \\
\bottomrule
\multicolumn{6}{@{}p{12.6cm}@{}}{\scriptsize The sole \mention{blast} token is coordinated \mention{Damn, and blast you to hell}. Shows can recur across heads. Token counts are verified lower bounds where candidates weren't displayed; the resulting share bounds are given in the text. The raw base isn't an audited count of construction opportunities.}
\end{tabular}
\end{table}

On those pronouns, 32 of 33 audited \mention{to hell} tokens (97\%) have \mention{damn} as head; the one exception is \mention{blast} coordinated with \mention{damn}, and neither \mention{darn} nor \mention{dang} contributes a token. In the \mention{for} V-ing frame \mention{damn} and its euphemisms supply only 11 of 48 tokens (23\%; \mention{damn} 10, \mention{darn} 1), against \mention{fuck} 17, \mention{screw} 13, \mention{bless} 6 and \mention{curse} 1, whose 37 tokens come from 36 shows.

Not every hit was displayed, so I bound both shares. On the shared pronouns, all 11 undisplayed \mention{to hell} hits have \mention{damn}, which can only raise its share (97--98\%). The 25 undisplayed \mention{for} V-ing hits include 2 with \mention{damn} and 23 with other heads, so the share of \mention{damn} and its euphemisms there lies between 15\% and 26\% however those lines turn out. The gap remains above 70 percentage points.

Conditioning on each head's raw shared-pronoun base leaves the contrast intact. \mention{Damn} has both PPs; \mention{screw}, \mention{fuck} and \mention{bless} have \mention{for} V-ing but no \mention{to hell} across 13,908 strings. Adjusting those bases for the sampled proportion of ordinary verbs and other uses preserves the contrast (Supplementary Section~S6). The following examples illustrate reason modification across heads; the full set is in the supplement.

\ea\label{ex:t7b-examples}
\ea \mention{Fuck her for killing herself.} \hfill (TV Corpus, \emph{13 Reasons Why})
\ex \mention{Well, bless him for trying.} \hfill (TV Corpus, \emph{The Grand Tour})
\ex \mention{Darn you for raising us so well.} \hfill (TV Corpus, \emph{Good Luck Charlie})
\z
\z

Syntactically, (\ref{ex:t7b-examples}a) links target and reason: in \mention{Fuck her for killing herself}, \mention{her} controls the understood subject of \mention{killing}, which binds reflexive \mention{herself}. Adjunct control is also compatible with a discourse afterthought (\mention{Fuck her! For killing herself, of all things.}), so control alone doesn't decide between attachment inside the clause and modification of the utterance. Semantically, the PP gives the grounds for the speaker's attitude, and an interlocutor can contest them (\mention{A: Fuck you for lying! B: I didn't lie!}).

The searches find 37 subjectless \mention{fuck}/\mention{screw} + pronoun + \mention{for} V-ing tokens and four with retaliatory particle verbs (Section~\ref{sec:beyond}), but none after finite or infinitival \mention{fuck}/\mention{screw} + pronoun. An audited base of ordinary transitive uses would be needed to compare their rates; these counts establish where the frame is attested.

Lexical NPs and overt \mention{God} also occur (\ref{ex:f2-examples}); the latter is compatible with Position~A for \mention{bless}. Searches with \mention{all}/\mention{both} after the pronoun reproduce the PP split (Supplementary Section~S6).

\ea\label{ex:f2-examples}
\ea \mention{Curse your father for making me feed him.} \hfill (TV Corpus, \emph{Married with Children})
\ex \mention{\dots but God bless the kids for trying to save it.} \hfill (TV Corpus, \emph{The Middle})
\z
\z

\mention{To hell} also occurs with non-referential \mention{it}: 52 raw \mention{damn it to hell} strings and one \mention{dang it to hell}, against none with \mention{fuck it} or \mention{screw it}, whose bases predict four or five each at \mention{damn}'s rate. The restriction extends beyond reference to a consigned entity; these counts leave lexical and semantic explanations open.

A search for head + \mention{it for} V-ing (\texttt{D14}) returns no TV hits for \mention{damn}, \mention{fuck}, \mention{screw}, \mention{bless}, \mention{sod} or \mention{bugger}. It omits \mention{curse} and \mention{darn}, which also take reason PPs. The zero is informative mainly for \mention{damn it}, with over a hundred tokens expected at its shared-pronoun rate. These are raw-string benchmarks: referential uses and formulae weren't separately audited (Supplementary Section~S6).

D14 leaves open both a semantic requirement for a target held responsible and a conventional restriction on formulaic \mention{Damn it!}. On the judgment in Section~\ref{sec:subtypes}, its \mention{it} is non-referential, though a speaker, addressee or other participant could supply a target in the discourse. Distinguishing the explanations requires evidence about the uses' reference and targets independent of the PP.

An inanimate lexical-NP target is attested once in GloWbE (\mention{darn the rains for messing with our social scene}; Section~\ref{sec:minced}). How widely inanimate targets accept \mention{for} V-ing across heads, and whether they depend on personification (\mention{Damn this printer for jamming again!}), remain questions for judgment studies.

\subsection{Beyond the sample: \mentionhead{thank}, \mentionhead{shame on}, \mentionhead{silly me} and \mentionhead{yay}}
\label{sec:beyond}

\mention{For} V-ing occurs well beyond the sample. \textcite[101--103]{fillmore1969judging} describes \mention{for} phrases with \mention{criticize}, \mention{scold}, \mention{blame} and \mention{praise} as identifying the situation involved in a judgment. The comparisons here include expressions of judging, thanking and retaliating.

After \mention{thank you}, the PP is frequent (\mention{thank you for coming}); it also occurs with the praise and blame formulae \mention{shame on}, \mention{good for} and \mention{kudos to}, with \mention{silly me} (2 tokens) and with retaliatory particle verbs (\mention{he will fuck you up for hurting me}, \mention{screw me over for recruiting Bo}; 4 tokens).

In these corpora it doesn't occur after \mention{lucky}, \mention{poor}, \mention{yay} or \mention{boo}, though \mention{yay me} and \mention{boo me} occur as raw strings (32 and 21 times). Outside them, \mention{yay} takes both a reason PP (\mention{Yay You For Making The Thread!!!}, \mention{Yay me for catching up on 40 years worth of Doctor Who History}\dots; DISboards, 22 April and 8 May 2006) and a lexical NP: a blog post welcoming free prescriptions has \mention{Yay the NHS!} (Bookmole's Blog, 3 May 2012). The lexical-NP contrast in Sections~\ref{sec:tv-profile} and~\ref{sec:bnc2014} is confined to those corpora; the blog example leaves attachment and productivity unresolved.

\mention{Thank} has 15,929 raw TV strings of \mention{thank} + pronoun + \mention{for} V-ing (\texttt{C2a}; \mention{thank you for coming} 3,408), against 99 for the sampled heads (\texttt{T7b}). The \mention{thank} strings include imperative \mention{thank him for}\dots\ and haven't been audited. They document frequent occurrence outside the sample; estimating how often \mention{thank} takes the PP would require an audited base. Whether \mention{thank} lacks \mention{to hell} and overt \mention{God} hasn't been queried.

\mention{Thank you} also takes sender and sincerity \mention{from} phrases (\mention{from all of us}, \mention{from the bottom of my heart}) in far greater absolute numbers than the sampled heads; no denominator was audited (Supplementary Section~S7). The sender's role is compatible with a reduced performative but doesn't diagnose a subject. Outside the corpora, \mention{Fuck you from the bottom of my heart} appears in a Facebook message quoted in \emph{In the Interest of O.W.} (Iowa Court of Appeals, no.~24-0862, 9 January 2025, p.~9).

\subsection{Overt \mentionhead{God}}
\label{sec:god-split}

The Tier-1 distribution recurs at scale in the TV Corpus. Of 4,084 tokens of \mention{God} + HEAD + NP, 4,031 have \mention{damn} or \mention{bless}. Setting aside \mention{it}-formulae and routine \mention{God bless you}, \mention{bless} has 686 and \mention{damn} 140, against 8 for \mention{curse}, 1 for \mention{fuck} (\mention{God fuck you}), and none for the other nine heads.

The minced heads take \mention{God} only in fixed \mention{it}-formulae. All 39 \mention{God dang/darn} tokens are on \mention{it}, and so are all 135 \mention{God/Gosh/Dog + darn/dang} strings in a follow-up search limited to pronouns (\mention{gosh darn it}, \mention{god dang it}).

Within \mention{bless}, overt \mention{God} depends on the expression. Across eight frequent strings (Supplementary Table~S2), \mention{God} precedes 80 of 86 tokens of four civic formulae (\mention{the United States}, \mention{the King}) but only 12 of 272 tokens of four idioms (\mention{your heart}, \mention{my soul}). The strings are illustrative rather than a sample, and the counts without \mention{God} are raw (they include song lyrics and surprise \mention{bless my soul}). Even so, 93\% against 4\% points to conditioning by expression rather than a uniform rule; an optative account could allow omission conditioned by formula.

\subsection{The bare-subjunctive \mentionhead{be} + participle frame}
\label{sec:passive}

Outside the active frame, some heads occur as participles after bare subjunctive \mention{be} (\ref{ex:diagnostics}f), and the form of the frame depends on the head. With \mention{damn} and \mention{bugger} it's a dismissal with the NP first. A search for NP + \mention{be damned/buggered} verifies 88 dismissals in 100 audited lines, across 73 shows (85 with \mention{damned}, as in \mention{The law be damned!}, and 3 with \mention{buggered}). \mention{Damned be} occurs only in four quotations or echoes of a line from \emph{Macbeth} (\mention{And damned be him that first cries}\dots).

With \mention{bless} the order is reversed and the force is a benediction. None of 18 lines of NP + \mention{be blessed} is a bare subjunctive (15 have \mention{may}, as in \mention{May your days be blessed}), whereas \mention{Blessed be} NP gives 83 benedictions. Most are liturgical or show-specific formulae (\mention{Blessed be the name of the Lord}; \emph{The Handmaid's Tale}'s \mention{Blessed be the fruit}), and formulaic \mention{blessed} also admits an adjectival analysis alongside a verbal-passive one; even so, parodies show the inverted frame taking new NPs (\mention{Blessed be the regulations}), contrasting sharply with uninverted NP + \mention{be damned}.

Secular \mention{fuck} and \mention{screw} and the minced forms barely have it. NP + \mention{be fucked/screwed} gives one dismissal (\mention{Memory be fucked}, \emph{True Detective}). The minced forms have \mention{be darned} almost only as a formula: in a 100-line sample, 99 are \mention{I'll be darned!} or \mention{be darned if}, and the other is the one dismissal that a search for NP + \mention{be darned/danged} also finds (\mention{Rules be darned, Richard.}, \emph{Call Me Fitz}).

\subsection{Euphemisms}
\label{sec:minced}

The minced forms show \term{partial euphemistic correspondence}: they occur in many of their sources' frames, in different proportions. Supplementary Table~S6 gives the full diagnostics. In GloWbE, 36 of 51 \mention{darn} + full-NP candidates are imprecatives, against 5 of 20 for \mention{dang}; most remaining \mention{dang} lines precede a clause (\mention{dang that list is long!}). Four \mention{darn} examples take \mention{for} V-ing, including \mention{darn the rains for messing with our social scene}.

Pronoun targets show a related difference. \mention{Darn} has nine targeted curses and one bleached oath in 13 lines; \mention{dang}'s two subjectless tokens in 11 lines are both bleached oaths (\mention{Well dang me sideways}). The samples are too small to show that targeted pronoun curses are unavailable for \mention{dang}. Zeros in finite and infinitival curse uses, and with \mention{to hell}, are similarly weak (Supplementary Section~S8).

Outside the pattern, the minced forms share attributive and adjective-modifying uses with \mention{damn} (\mention{that darn donkey}, \mention{pretty darn good}); \mention{dang} modifies the adjective in 90 of 100 sampled lines. They also share the finite oath uses described in Section~\ref{sec:passive}. These correspondences extend beyond the sampled pattern and can be stated in lexical entries or construction networks.

The six euphemisms of \mention{fuck} give 68 imprecatives in 100 sampled TV lines from 223 hits. \mention{Frak} supplies 40, all but one from \emph{Battlestar Galactica} and \emph{Caprica}. Alongside open NPs (\mention{Frig the dip}, \emph{The Sopranos}), \mention{eff} and \mention{frak} occur outside the pattern (\mention{don't eff this up}, \mention{where the frak}). Subtitle distributions show franchise-specific usage and may also reflect broadcast censorship.

\ea\label{ex:f3-examples}
\ea \mention{You know what? Eff Christmas. Eff you.} \hfill (TV Corpus, \emph{Saturday Night Live})
\ex \mention{Eff you both! Eff you all to H!} \hfill (TV Corpus, \emph{Call Me Fitz})
\z
\z

\mention{Eff Christmas} has the open NP expected of the pattern. \mention{Eff you all to H!} combines a euphemism of \mention{fuck} with the extension found in \mention{damn you all to hell}, consistent with analogy across formulae. The extended examples and audits are in Supplementary Section~S8.

\subsection{Unscripted speech: Spoken BNC2014}
\label{sec:bnc2014}

Spoken BNC2014 supplies an unscripted comparison (Supplementary Section~S9). The six living heads dominate, while only four of 22 archaic/minced candidates are imprecatives. Among 100 inspected determinative/genitive lines, 48 are subjectless tokens, including five open full NPs (\mention{Fuck this boat club}, \mention{bless that poor girl who was acting}). No expressive lexical-NP complement occurs in 17 \mention{yay}/\mention{boo} candidates.

Overt \mention{God} again occurs only in theological optatives: ten \mention{god damn it} and five \mention{god bless}. The PP searches find one \mention{damn it to hell} and three \mention{for} V-ing tokens in three conversations (\mention{bless him for carrying}, \mention{damn you for being busy}, \mention{bless him for coming on TV}). These counts are compatible with the TV-rate benchmarks of about three and two to three respectively.

No imprecative occurs among eleven \mention{fuck}-euphemism candidates. But removing the dominant \emph{Battlestar Galactica}/\emph{Caprica} franchise from the TV calculation lowers the expected speech count from about five to two, so the zero says little. The detailed raw-string rate comparisons and their limitations are in Supplementary Section~S9.

% ===========================================================================
\section{Discussion: Understood subjects, meaning, scope, and categories}
\label{sec:discussion}
% ===========================================================================

The discussion takes four questions in turn: which syntactic analyses the binding and distributional facts constrain (Section~\ref{sec:ruled-out}), what the same facts require of semantic accounts (Section~\ref{sec:meaning}), which generalizations the observed material establishes, and which explanations of them remain candidates (Section~\ref{sec:convergence}).

\subsection{What the diagnostics exclude, and the head slot}
\label{sec:ruled-out}

The results constrain the accounts in Section~\ref{sec:four-positions}. Non-reflexive \mention{you} excludes an understood addressee subject for the \mention{you} forms, though a speaker subject remains possible (\mention{I damn you}). Non-reflexive \mention{me} in self-curses (\mention{Damn me for agreeing to it}) excludes a speaker subject, though with someone present an addressee subject remains possible (the imperative reading, with an understood \mention{you}). So no single addressee or speaker subject fits all the sampled heads. These tests concern a syntactic subject, not whose attitude the expression conveys. In \mention{God damn you} the subject is God and the attitude is the speaker's; \textcite[173--175]{potts2007} treats the holder of an expressive attitude as a contextual judge, by default the speaker, and \textcite{harrispotts2009} find that readings of appositives and epithets that aren't the speaker's are rare in use but systematic. So speaker orientation is no evidence of a speaker subject.

A mixed analysis, with a speaker subject for the \mention{you} forms and an addressee subject for the \mention{me} forms, comes under further pressure wherever self-curses occur with no one present, since binding still excludes the speaker and, with no one present, there's no conversational addressee for an understood subject to refer to. First-person narration supplies such cases. A narrator who has just \enquote{found myself alone} continues \mention{Tricked. Damn him. And damn me for falling for it.} (COCA, Kelley Armstrong, \emph{Broken}, 2006), and \mention{Damn me for shooting my mouth off that night} (COCA, \emph{Fantasy \& Science Fiction}, 2007) is addressed to no one in the scene. To my ear, \mention{Fuck me for being an idiot!}, muttered alone, is possible, and bleached \mention{Fuck me!} occurs alone too (on dropping something heavy, say), though its \mention{me} may not be referential and so says nothing about a speaker subject.

Solitary self-curses don't rule out every analysis with an understood subject. Solitude doesn't rule out apostrophe to an absent addressee (compare imperative wishes such as \mention{God, help me!}, where \mention{God} is a vocative and the understood imperative subject is second person, distinct from optative \mention{God help me!}, where \mention{God} is a third-person subject), and scarce overt \mention{God} on secular \mention{fuck} (a single TV \mention{God fuck you}, Section~\ref{sec:god-split}) doesn't prove that an understood third-person subject is unavailable. Even so, both alternatives need support. An understood subject referring to an absent addressee needs independent motivation where no one is addressed. For God, \textcite[6]{quang1971} already applied the binding test: \mention{Damn God} is acceptable and \ungram{\mention{Damn Himself}} isn't, and the same holds for \mention{goddamn} and \mention{fuck}, so a \mention{God} subject restored throughout would wrongly predict the reflexive wherever God is the target. The 140 non-\mention{it} tokens of overt \mention{God damn} show a conventional relation between \mention{damn} and \mention{God}, not that every bare \mention{damn} derives from it. For \mention{fuck}, then, the solitary cases favour an analysis with no understood subject (Positions~B--D), if that judgment holds; for \mention{damn}, they leave room only for a \mention{God} subject limited to some uses. The secure exclusion is that no single understood subject, whether addressee, speaker or \mention{God}, accounts for all the uses.

Open lexical NPs separate these heads from the six primary interjections searched. Other interjection-like heads take dependents (\mention{shame on you}, \mention{yay me}), so a complement-taking interjection phrase (Position~B) remains open. Finally, \mention{confound}, \mention{blow}, \mention{stuff} and the euphemisms show which words fill the head slot \citep{mcenery2006,napolihoeksema2009,ljung2011}. Whether a head is a recent recruit or a survival is a diachronic question these counts can't answer, though \mention{confound}'s concentration in period drama suggests survival.

\subsection{What the distributions require of semantic accounts}
\label{sec:meaning}

Semantic accounts distinguish expressive meanings that need an argument from those that don't. \textcite[175]{potts2007} states the meaning of attributive \mention{damn} for an utterance with an entity as its semantic argument, and \textcite[6, 27]{gutzmann2013} calls argument-seeking items such as attributive \mention{damn} functional and saturated ones such as the interjection \mention{Damn!} isolated. An expressive attitude can be directed at an individual, at \enquote{some specific feature of the current state of affairs}, or be \enquote{just general, undirected emotion} \citep[173]{potts2007}. For the subjectless forms themselves, \textcite[10--11]{portner2026} proposes a social-relation meaning that varies with the argument: \mention{Damn it} has the meaning of \mention{Damn you} \enquote{applied to the general situation}, \mention{Fuck it} \enquote{has to do with action}, and \mention{Fuck me} is mirative. Three of the distributions bear on these accounts.

First, \mention{damn}, \mention{fuck}, \mention{screw} and \mention{bless} take reason PPs with other complements but have no \mention{it for} V-ing in \texttt{D14} (Section~\ref{sec:dissoc}). On Portner's account, \mention{Damn it} concerns the general situation and \mention{Fuck it} a course of action. Their contrast with referent-directed uses suggests that target type may matter to reason modification. D14 establishes an attestation contrast by pronoun; it leaves the target requirement and the formula-specific alternative open.

Second, first-person targets occur. \textcite[15]{portner2026} counts it an advantage of the social-relation analysis that it \enquote{explains the fact that imprecative quasi-verbs cannot take first person arguments}, since a speaker has no social relation to themself, and he marks \mention{Damn me} as impossible. The TV Corpus has seven strings of a head + \mention{me} + \mention{for} V-ing, after four heads. All five displayed lines are self-directed curses: three are sincere (\mention{And damn me for agreeing to it}, \emph{Babylon 5}; \mention{fuck me for ordering it in the first place}, \emph{Hand of God}; \mention{Fuck me for giving a shit}, \emph{Shameless}) and two are sarcastic (\mention{Well, fuck me for ruining your night}, \emph{Dexter}; \mention{screw me for thinking you might be open to changing your mind}, \emph{Jericho}), and COCA has two narrated self-curses addressed to no one (Section~\ref{sec:ruled-out}). These are targeted readings, not the mirative one. Each has a reason PP stating what the speaker is blamed for, and the samples don't show whether bare \mention{Damn me!} occurs as a self-curse without one. The counts are small, but an account has to cover these uses, whether with a relation of the speaker to themself, a separate meaning, or a revised explanation.

Third, the reason PP isn't confined to expressive meaning. It occurs with \mention{thank you}, \mention{shame on you} and \mention{kudos to} (Section~\ref{sec:beyond}), and \mention{I thank you for coming} has the same PP with a described act of thanking. Whatever meaning the imprecatives have, the PP's distribution follows acts of praising, blaming, thanking and retaliating (\mention{he will fuck you up for hurting me}), whether those acts are expressed or described.

A hearer can deny the grounds directly (\mention{A: Fuck you for lying! B: I didn't lie!}; Section~\ref{sec:dissoc}), so their content can be challenged. Whether \mention{Fuck her for killing herself} asserts or presupposes that she killed herself, and whose attitude a reported or embedded curse conveys, are questions for judgment studies of the kind \textcite{harrispotts2009} ran on appositives and epithets.

\subsection{The scope of each generalization}
\label{sec:convergence}

The diagnostics give an uneven, head-specific profile (Table~\ref{tab:retention}).

\begin{table}[tbp]
\centering
\caption{Properties of each head and its associated expressions (bare HEAD + NP, overt-\mention{God} optatives, \mention{be} + participle frames), from audited Tier-1 and Tier-2 diagnostics. A row gathers forms related to one head without implying that they belong to one lexical entry. TV PP counts use the shared pronouns \mention{you/him/her/them} unless marked \mention{it}-only. Bracketed figures beside zeros in the PP columns give the raw TV pronoun base (\texttt{T2} \mention{you} + \texttt{T3} \mention{him/her/them}).}
\label{tab:retention}
\footnotesize
\setlength{\tabcolsep}{3.2pt}
\begin{tabular}{@{}>{\raggedright\arraybackslash}p{1.85cm}>{\raggedright\arraybackslash}p{1.9cm}>{\raggedright\arraybackslash}p{1.85cm}>{\raggedright\arraybackslash}p{1.7cm}>{\raggedright\arraybackslash}p{1.65cm}>{\raggedright\arraybackslash}p{2.05cm}>{\raggedright\arraybackslash}p{1.1cm}@{}}
\toprule
Heads (stratum) & Open full NP & Non-\mention{it} \mention{God} & Goal \mention{to hell} & Reason \mention{for} V-ing & Bare-subj.\ \mention{be}+part. & Coord. \\
\midrule
\mention{damn} (theol.) & Tier 1, TV, BNC14 & 140 (TV) & 32/33 (TV); 1 (BNC14) & 10 (TV); 1 (BNC14) & Uninverted: 85 of 100 audited; \mention{Damned be} only in 4 \emph{Macbeth} quotations (TV) & 11 (\mention{blast}) \\
\mention{bless} (theol.) & Tier 1, TV, BNC14 & 686 (TV); 5 (BNC14) & 0 [3,113] & 6 (TV); 2 (BNC14) & Inverted only: 83 \mention{Blessed be} NP; 0 of 18 bare \mention{NP be blessed} (TV) & 0 \\
\mention{fuck}/\mention{screw} (secular) & Tier 1, TV, BNC14 & 1 (TV, \mention{fuck}); 0 (BNC14) & 0 [10,795] & 30 (TV), plus 3 with \mention{all} & 1 (TV, \mention{Memory be fucked}) & 0 \\
\mention{sod}/\mention{bugger} (secular UK) & Tier 1, TV & 0 & 0 [75] & 0 [75] & 3 audited (4 non-OCR \mention{buggered}) & 0 \\
\mention{darn}/\mention{dang} (minced) & TV; GloWbE \mention{darn} 36 of 51, \mention{dang} 5 of 20 & 0 (39/39 \mention{it}) & \mention{it}-only (1 \mention{dang it}) & 1 (TV, \mention{darn}); 4 (GloWbE, \mention{darn}) & 1 \mention{Rules be darned}; finite oath \mention{I'll be darned!} (99/100) & 0 \\
\bottomrule
\end{tabular}
\end{table}

Even the two theological heads differ: \mention{damn} concentrates \mention{to hell}, while \mention{bless} has civic \mention{God} formulae, and their bare-subjunctive participles occur in opposite orders with opposite force. The broader profile combines properties shared by most heads, restrictions on particular expressions, and correspondences extending beyond the pattern. A common historical stratum leaves differences to explain.

Each kind of observation licenses some inferences and not others (Table~\ref{tab:inferences}).

\begin{table}[tbp]
\centering
\caption{What each kind of observation supports, and what it doesn't establish by itself.}
\label{tab:inferences}
\small
\begin{tabular}{@{}>{\raggedright\arraybackslash}p{3.3cm}>{\raggedright\arraybackslash}p{4.5cm}>{\raggedright\arraybackslash}p{4.5cm}@{}}
\toprule
Observation & Supports & Doesn't establish by itself \\
\midrule
Open lexical NPs & expansion beyond fixed pronoun formulae & a single lexical category or representation \\
Non-reflexive \mention{you} and referential \mention{me} & limits on understood addressee and speaker subjects & a lexical category, or that no subject is understood \\
Concentration of \mention{to hell} & strong conditioning by \mention{damn} and expressions related to it & that a historical lineage is the relevant unit \\
Reason PPs across forms & a recurring relation between a target, an attitude and its grounds & one attachment structure, or a semantic licensing condition \\
No \mention{it for} V-ing for six queried heads (\texttt{D14}) & an attestation contrast with other pronoun complements for four heads & a requirement for a referential or responsible target \\
Self-directed \mention{me} with a reason PP & first-person targets in the pattern & that they need the PP, or share the second-person meaning \\
Euphemisms in their sources' frames & relations between words that extend beyond the pattern & that a frame is unavailable where a rare form wasn't found \\
\bottomrule
\end{tabular}
\end{table}

Lexical, historical, semantic, constructional and pragmatic classifications of these expressions each carry part of this information (the pragmatic categories of Section~\ref{sec:intro} cut across the rest, since the same form curses, blesses, thanks and dismisses), and none has priority for every question.

The semantic proposal cuts within heads: formulaic \mention{Damn it!} is situation-directed, while \mention{Damn him!} can evaluate a referent. Whether that distinction, identified independently of the PP, predicts reason modification across heads remains open (Section~\ref{sec:meaning}).

A lineage-sensitive account of \mention{to hell}, an evaluative relation for \mention{for} V-ing and a productive pattern for open NPs are candidate explanations with different standing: the first rests on one head, the second partly on how the uses were classified (Section~\ref{sec:subtypes}), and the third doesn't decide between one construction and lexical entries with similar requirements. Because most of the comparison is between heads that a verbal analysis would put in one lexical category, the results measure variation within that category, not how much it predicts compared with others.

This graded profile defines what any analysis has to capture, while these results still underdetermine the choice among the three accounts compared here: a complement-retaining interjection phrase \citep[1361 n.\ 70]{huddleston2002}, defective verbal entries formulated in a lexicalist framework \citep{mullerwechsler2014}, and a structured usage-based constructional network \citep{goldberg2003,goldberg2019,croft2001,bybee2010,diessel2019,schmid2020}, paralleling sign-based continua \citep{sag2012sign}. That's partly because, as Section~\ref{sec:four-positions} notes, they answer different questions. Both enriched lexical entries and anchored construction networks can state head-specific properties, default accusative case \citep{schutze2001}, evaluative \mention{for} V-ing, and horizontal euphemistic correspondences. Testing acceptability across \mention{yay}/\mention{boo}, preposition-headed dismissals (\mention{to hell with} X, \mention{down with} X), and the productivity and personification sensitivity of inanimate \mention{for} V-ing remains a productive experimental agenda.

% ===========================================================================
\section{Limitations and conclusion}
\label{sec:conclusion}
% ===========================================================================

A matched American spoken benchmark remains for future work, and a single coder adjudicated the lines. Corpus counts establish attestation; judgment and audio studies would be needed for attachment, graded acceptability and the role of personification in inanimate reason-PP targets. The results concern these heads in this form, rather than cursing and blessing generally.

Within that scope, generalizing from \mention{damn} or \mention{fuck} obscures differences among heads and uses. \mention{To hell} concentrates on \mention{damn}; reason modification reaches beyond the sampled heads to \mention{thank you} and other evaluative expressions; euphemisms occur in many of their sources' frames in different proportions. Non-reflexive complements constrain understood subjects, and self-curses require semantic accounts to accommodate first-person targets.

These findings bear on categorization through different questions. Lexical, historical, semantic, constructional and pragmatic distinctions each contribute information, but the inherited set of heads is a sampling frame, not a category selected by the properties studied. Lineage-sensitive retention, an evaluative relation licensing reason modification and a productive pattern for open NPs remain candidate explanations with different evidential standing. The results constrain uniform accounts of subjects and extensions while leaving complement-taking interjection phrases, defective lexical entries and lexically specified constructions available as representations.

% ===========================================================================
\section*{Data availability and disclosures}
\ifblind
Extraction scripts, line-level Tier-1 adjudication tables, parsed Tier-2 frequency summaries, and all exact query definitions (Supplementary Table~S1) will be archived in an anonymized repository for review and publicly upon acceptance, excluding restricted third-party corpus texts. An AI assistant was used for Python extraction and bookkeeping scripts, for manuscript drafting, and for first-pass line codes on some concordance samples; the author designed the study, ran every web-corpus query by hand, checked every line-level code, audited all analytical claims, and takes full responsibility for the text.
\else
Extraction scripts, line-level Tier-1 adjudication tables, aggregated Tier-2 summaries, and all exact query definitions (Supplementary Table~S1) are archived publicly at \url{https://github.com/BrettRey/english-imprecatives} (excluding restricted third-party corpus texts). An AI assistant was used for Python extraction and bookkeeping scripts, for manuscript drafting, and for first-pass line codes on some concordance samples; the author designed the study, ran every web-corpus query by hand, checked every line-level code, audited all analytical claims, and takes full responsibility for the text.
\fi

\clearpage
\printbibliography

\end{document}
